% Requiere en el preámbulo:
% \usepackage{hyperref}
% \usepackage{xurl} % opcional, ayuda a cortar URLs largas

\section{REPOSITORIO DIGITAL DEL INFORME}

En la presente hoja se consigna la dirección URL del repositorio digital en Google Drive donde se encuentra la documentación complementaria del informe. Además, se resume el tipo de materiales que respaldan la trazabilidad documental del plan.

\vspace{0.5cm}

\noindent\textbf{Dirección URL del repositorio:}

\vspace{0.3cm}

\noindent
\url{https://drive.google.com/drive/folders/1t_vPJ7YbWgOVJ2CKJ4A22lGIAcpq6uH5?usp=drive_link}

\vspace{0.8cm}

\noindent\textbf{Contenido del repositorio:}

\begin{enumerate}
    \item Papers en formato PDF y sus respectivos resúmenes, con los archivos renombrados según el título del paper.
    
    \item Ejemplos de datos compilados o parametrizados para los casos de estudio.
    
    \item Normas, estándares, guías y manuales relacionados con el tema de investigación.
    
    \item Archivos originales de los documentos entregados en PDF y \LaTeX.
    
    \item Archivo Excel que contiene la matriz de consistencia.
\end{enumerate}

\vfill

\noindent
Este repositorio constituye el respaldo digital de los documentos, fuentes, datos y anexos utilizados para el desarrollo del informe. Por tanto, apoya la trazabilidad de las fuentes y de los materiales complementarios del plan.




